% sec_conclusion.tex — Section 6: Conclusion (compressed)
% ICGNC 2026, Springer svproc template

\section{Conclusion}
\label{sec:conclusion}

This paper presented TAGA-LLM, a two-stage framework that bridges MARL cooperative
performance with LLM interpretability for multi-UAV mission planning.  The TAGA
mechanism introduces three learnable, physics-informed attention
factors---distance decay $f_d$, task awareness $f_m$, and information freshness
$f_t$---that model realistic communication constraints beyond static-graph
assumptions.  The MARL-to-LLM distillation pipeline, combining topology-aware
prompt encoding, chain-of-thought supervision, and counterfactual augmentation,
transfers cooperative strategies into a Qwen-4B model that provides transparent
reasoning, instruction controllability under the five mission directives
(\textsc{bal}--\textsc{coo}), and efficient edge deployment
(${\sim}70$\,ms/step with 4-bit quantization).

\paragraph{Limitations.}
The current framework assumes homogeneous UAVs and a simplified grid-based
environment.  The distillation quality depends on the coverage of the expert
trajectory dataset, and the 256-token prompt truncation may lose information in
dense scenarios.  Additionally, the chain-of-thought supervision relies on
template-based rationale generation rather than human annotations.

\paragraph{Future work.}
Promising directions include extending to heterogeneous swarms with
role-specific policies, incorporating online adaptation for non-stationary
environments, and validating on real-world UAV testbeds with physical
communication constraints.
